\section{A human-readable conditional main proof}
\label{sec:conditional-main-v42}

The argument below reaches RH and GRH for Dirichlet $L$-functions
from two explicitly stated arithmetic inputs.  It gives the whole
terminal proof in ordinary mathematical notation, including the
normalizations, the dependence of the constants and the passage from
a one-sided prime-sum estimate to zeros.  The two arithmetic inputs
are not asserted here.  The current reconstruction has proved many
of their component implications, but their common construction and
the five quantitative estimates are still incomplete.  The final
table locates every v0.41 obligation in this main argument.  The
adjoint theorem in the next section closes one reference-matching
obligation; it does not supply these two arithmetic inputs.

\subsection{The fixed analytic target}

Let $\Lambda$ be the von Mangoldt function and let $\gamma$ be
Euler's constant.  For $y>0$ put
\begin{align*}
 G(y)&=\sum_{r\le y}\Lambda(r)(y/r-1)
                         -y\log y+(1+\gamma)y,\\
 G_\chi(y)&=\sum_{r\le y}\Lambda(r)\chi(r)(y/r-1)
                         +yL'/L(1,\chi)\qquad(\chi\ne1).
\end{align*}
Here $\chi$ is primitive, and $G_1$ means $G$; the sum is empty
if $y<1$.  The plus sign before $yL'/L(1,\chi)$ cancels the
unwanted linear Mellin pole.  Fix a primitive target $\psi$ of
order $o$.  Let $\mathcal C_\psi$ be the set of distinct primitive
inducers of its powers.  Its positive signal is
\[
 H_\psi(x)=G(x)+\frac1{3o+1}
             \sum_{\chi\in\mathcal C_\psi\setminus\{1\}}
                                      \operatorname{Re}G_\chi(x).
 \tag{CMP.1}
\]
This family and these coefficients are fixed throughout the proof.
For $\psi=1$ it is simply $G$.

The standard analytic facts used below are continuation and the
functional equations of primitive Dirichlet $L$-functions, their
Euler products, and their nonvanishing on $\operatorname{Re}s=1$.
The reconstruction already supplies the needed nonvanishing and
primitive-character transitions in QW14 and GH6.  Complex analysis
then gives CMP4 below.  No unproved zero-free half-plane inside the
critical strip is among the analytic inputs.

\subsection{First arithmetic input: one compatible finite construction}

At a scale $x\ge3$ and an accuracy parameter $\varepsilon>0$,
the following objects and assertions constitute a
\emph{compatible finite datum}.  They are deliberately stated on
one set of objects.

Take a finite abelian group $A$ of order $n$, a subgroup $A_0$ of
index $o$, and write
\[
 G_0=(\mathbb Z/2\mathbb Z)^2\times A,\quad
 \Phi=\{(0,b,h):b\in\mathbb Z/2\mathbb Z,\ h\in A\},
 \quad f=|\Phi|=2n.
\]
Set $c_0=(1,0,0)$, $c_\infty=(1,1,0)$,
$\chi_d(a,b,h)=(-1)^a$ and $\chi_q(a,b,h)=(-1)^b$.
Supply a positive modulus $N_0$ and an onto homomorphism
\[
 r:(\mathbb Z/N_0\mathbb Z)^\times\longrightarrow G_0.
\]
For every character $\xi$ of $G_0$, let $\xi^*$ be the primitive
inducer of $\xi\circ r$, and set
$\mathcal G_\xi=G_{\xi^*}$ and
$\ell_\xi=L'/L(1,\xi^*)$ for $\xi\ne1$.
These are actual equalities, not freely chosen approximations to
the logarithmic derivatives.

For $p\mid N_0$, define $I_p$ to be the image of the $p$-primary
CRT unit factor.  Define $F_p$ modulo $I_p$ by the class of $p$
in the prime-to-$p$ factor.  For $p\nmid N_0$, put $I_p=0$ and
$F_p=r(p)$.  The choice of a representative of $F_p$ will not
matter.  As in the source host, supply a finite prime set $S$ and
a prime $d\notin S$ such that
\[
 I_p=0\ (p\notin S\cup\{d\}),\qquad
 I_p\subset\Phi, F_p\in\Phi\ (p\in S),\qquad
 I_d=\{0,c_0\}.                                    \tag{CMP.2}
\]

Put $\mathcal B=\{\xi:\xi(0,0,A_0)=1\}$; it has $4o$
elements.  Supply a subset $K\subset\mathcal B$ containing $1$
such that $\xi(c_0)=1$ for $\xi\in K$, $\chi_q\notin K$, and
$\xi\mapsto\xi^*$ identifies $K$ with
$\mathcal C_\psi$.  Supply nonempty $T\subset A$ and
$T_2\subset A_0$, and set
\[
 n_T=|T|,\quad n_2=|T_2|,\quad d_1=n_T/n,\quad d_2=n_2/n.
\]
The required bias in every coset is
\[
 \left|\#(T\cap(h+A_0))-
       \frac{3+\mathbf1_{h\in A_0}}{3o+1}\,n_T\right|
                   \le x^{-1+\varepsilon}n_T.
 \tag{CMP.3}
\]
For constants $K_0,C_E$ that stay fixed as the windows move,
require
\[
 |\ell_\xi|\le K_0\ (\xi\in K\setminus\{1\}),\qquad
 |\ell_\xi|\le C_E\log x\ (\xi\in\mathcal B\setminus K),
 \quad
 |\mathcal G_\xi(x)|\le x^{1/2+\varepsilon}
                 \ (\xi\in\mathcal B\setminus K).
 \tag{CMP.4}
\]

Choose $m\ge1$ integral, $b,b',X',k>0$, $\eta,\theta_0\ge0$,
$0\le a_E\le2xb$, numbers $a_t>0$ for $t\in T$, and positive
weights $w_{u,j}:\Phi\to\mathbb R_{>0}$ for
$u\in A_0$, $1\le j\le m$.  The thick and thin weights satisfy
\begin{gather*}
 |a_t/b-x|\le1,\qquad b'X'\le xb,\\
 b\le w_{u,j}(0,1,s+u)\le b+\eta+1\quad(s\in T_2),\\
 w_{u,j}(\rho)\le2b'
   \quad\bigl(\rho\notin\{(0,1,s+u):s\in T_2\}\bigr),\\
 |\mathcal G_\xi(y)|\le y\theta_0
           \quad(\xi\in\widehat G_0, y\ge X'/3),\\
 2f(x+1)\theta_0\le x^{-1},\qquad
 6f(x+1)(\eta+1)\le b/x.                         \tag{CMP.5}
\end{gather*}
The scalar sizes are
\[
 d_1\ge\tfrac12x^{-\varepsilon/8},\quad
 x^{-\varepsilon}\le d_2\le1,\quad
 x^{1/2-\varepsilon}\le m\le x^{1/2+\varepsilon/2},\quad
 xm\le k\le n^2x^{1/2-\varepsilon}.                \tag{CMP.6}
\]
Define the positive normalizations
\[
 Q=n_Tmb,\qquad D=Qd_2/2,
\]
and the error budget
\[
 B=n_Ta_E(1+m^2/k)
       +m\left(2nmbx^\varepsilon+
            \frac{n_2kbx^\varepsilon}{n(nx^{-\varepsilon})}\right).
 \tag{CMP.7}
\]
The symbols $D$ and $B$ have this meaning only in this main proof;
in particular $D$ is a scalar, not a replacement geometric object.

\subsection{Second arithmetic input: a uniform upper bound for one mean gap}

For $t\in A$, the bijection
$\rho_t(g)=c_\infty+g+(0,0,t)$ maps $G_0\setminus\Phi$ to
$\Phi$.  Put $(z)_+=\max(z,0)$ and define
\begin{align*}
 \mathcal A(a,t,w)
 &=\sum_{p}\sum_{j\ge1}\log p\,
       \frac2{|I_p|}
       \sum_{\substack{g\notin\Phi\\g-jF_p\in I_p}}
                    (a/p^j-w(\rho_t(g)))_+,\\
 \mathcal M(a,w)
 &=a\left(\log a-1-\gamma+\ell_{\chi_d}
                         -\frac1f\sum_{\rho\in\Phi}\log w(\rho)\right),\\
 \Gamma_u
 &=\sum_{t\in T}\sum_{j=1}^m
       \{\mathcal A(a_t,t,w_{u,j})
                         -\operatorname{Re}\mathcal M(a_t,w_{u,j})\},\\
 \overline\Gamma&=\frac1{|A_0|}\sum_{u\in A_0}\Gamma_u.
                                                               \tag{CMP.8}
\end{align*}
All prime-power sums here are finite: use a cutoff above all ratios
$a_t/w_{u,j}(\rho)$.  Conditions (CMP.2) make this precisely the
source's asked prime sum: primes in $S$ give zero, each unramified
term outside $\Phi$ has coefficient $2$, and the exceptional prime
$d$ has coefficient $1$ on its inertia coset outside $\Phi$.
The $d$ term is retained even when a particular cutoff makes it zero.

The quantitative input is the single inequality
\[
                    \overline\Gamma\le x^{C\varepsilon}B.
 \tag{CMP.9}
\]
The source's pointwise inequalities $\Gamma_u\le x^{C\varepsilon}B$
for every $u$ imply (CMP.9).  The next proof needs only (CMP.9);
this is a genuine weakening of the sufficient arithmetic input.
It allows cancellation between different translates before any
upper estimate is applied.

The complete order of choices is as follows:
\begin{quote}
For each fixed primitive $\psi$, choose $C\ge0$ and
$\varepsilon_0>0$.  For each $0<\varepsilon<\varepsilon_0$,
choose $K_0,C_E$ and any further constants of the construction.
For every $Y_0$, construct $Y\ge\max(3,Y_0)$ such that, for
every real $x\in[Y,Y^{3/2}]$, one compatible datum satisfies
(CMP.2)--(CMP.6) and its own inequality (CMP.9).
\end{quote}
Equivalently, the first input produces a common family on
arbitrarily late windows, and the second is uniform on that same
family after one further threshold.  The family may depend on
$\varepsilon$; the datum may depend on $x$.  The target,
coefficient $C$, and positive character weights may not be chosen
anew for each window.  Existence of separate data satisfying
separate clauses is insufficient.

% BEGIN CMP1
\begin{lemma}[CMP1: primitive values and the exact balance identity]
\label{claim:CMP1}
Every compatible datum satisfies
\[
 \mathcal A(a,t,w)-\mathcal M(a,w)
   =\frac1f\sum_{g\notin\Phi}\sum_{\xi\in\widehat G_0}
       \overline{\xi(g)}w(\rho_t(g))
                      \mathcal G_\xi(a/w(\rho_t(g))).
 \tag{CMP.10}
\]
\end{lemma}
\begin{proof}
For every prime power,
\[
 \xi^*(p^j)=\frac1{|I_p|}\sum_{h\in I_p}\xi(jF_p+h).
\]
If $\xi$ is nontrivial on inertia, both sides are zero: its
primitive conductor is divisible by $p$.  Otherwise it factors
through the prime-to-$p$ modulus, and both sides are $\xi(jF_p)$.
Character orthogonality consequently gives
\[
 \frac1f\sum_\xi\overline{\xi(g)}\xi^*(p^j)
          =\frac2{|I_p|}\mathbf1_{g-jF_p\in I_p}.
\]
Substitute into (CMP.8) and interchange finite sums.  For a
nonprincipal character, replacing the uncentered prime sum by
$w\mathcal G_\xi(a/w)$ subtracts $a\ell_\xi$.
For the principal character the extra terms are
$a\log a-a\log w-(1+\gamma)a$.
The character sum over $G_0\setminus\Phi$ equals $f$ at $1$,
$-f$ at $\chi_d$, and zero elsewhere.  These terms give exactly
$\mathcal M$, including its plus sign before $a\ell_{\chi_d}$.
This proves (CMP.10) at ramified as well as unramified primes.
\end{proof}
% END CMP1

% BEGIN CMP2
\begin{lemma}[CMP2: average first, then keep the positive target]
\label{claim:CMP2}
For all sufficiently large $x$, with threshold depending only on
$o,K_0,C_E,\varepsilon$, the same datum satisfies
\[
 H_\psi(x)\le\frac{\overline\Gamma}{D}
        +(4+4o)x^{1/2+\varepsilon}+\frac4{d_2x}.
 \tag{CMP.11}
\]
This conclusion uses the structural and cutoff conditions but not
the gap bound (CMP.9).
\end{lemma}
\begin{proof}
We recall the elementary uniform estimates established in WB3.
There is an absolute $C_0$ such that, for $y,z\ge1$,
\begin{align*}
 |G(y)-G(z)|&\le C_0|y-z|,\\
 |\mathcal G_\xi(y)-\mathcal G_\xi(z)|
 &\le |y-z|\{\log\max(y,z)+C_0+|\ell_\xi|\}\quad(\xi\ne1),\\
 |w\mathcal G_\xi(a/w)-w'\mathcal G_\xi(a/w')|
 &\le3a|w-w'|/\min(w,w').
                                                               \tag{CMP.12}
\end{align*}
For completeness the only arithmetic estimates in them are
$\sum_{r\le y}\Lambda(r)\le3y$ and
$\log k-1\le\sum_{r\le k}\Lambda(r)/r\le\log k+3$.
The first follows from central binomial coefficients and summing
prime powers; the second follows from
$\log(k!)=\sum_{r\le k}\Lambda(r)\lfloor k/r\rfloor$.
Differentiate the ramp sums between integers for the first two
inequalities in (CMP.12).  For the last use the Lipschitz property
of $(a/r-w)_+$; the nonprincipal linear term cancels exactly.
For the principal character the additional logarithmic difference
is bounded by $a|w-w'|/\min(w,w')$, and the difference of the two
nonnegative terms is still bounded by the displayed constant $3$.

Replace each thick weight by $b$ in (CMP.10).  The error per
character is at most $3(x+1)(\eta+1)$ by (CMP.12).
At a thin point,
$a_t/w\ge(x-1)b/(2b')\ge X'/3$, so its whole term has modulus
at most $(x+1)b\theta_0$.  There are $f$ points and $2f$
characters.  After division by $f$, (CMP.5) bounds the total error
for each $t,j$ by $2b/x$.
The ideal thick term is
\[
 \mathcal T_u(t,y)=\frac b f\sum_{s\in T_2}\sum_\xi
             \xi(-(1,0,s+u-t))\mathcal G_\xi(y).
\]
Average this term over $u\in A_0$.  The characters outside
$\mathcal B$ vanish, whereas those in $\mathcal B$ are trivial
on both $u$ and $s$.  Since $f=2n$, the result is
\[
 \frac1{|A_0|}\sum_u\mathcal T_u(t,y)
       =\frac{bd_2}{2}\sum_{\xi\in\mathcal B}
                                   \xi(1,0,t)\mathcal G_\xi(y).
\]
Thus (CMP.10), before any upper bound on $\Gamma_u$, gives
\[
 \frac{mbd_2}{2}\sum_{t\in T}
       \operatorname{Re}\sum_{\xi\in\mathcal B}
               \xi(1,0,t)\mathcal G_\xi(a_t/b)
                  \le\overline\Gamma+2Q/x.          \tag{CMP.13}
\]

For a nontrivial character of $A/A_0$, (CMP.3) makes its mean on
$T$ differ from $c=(3o+1)^{-1}$ by at most
$ox^{-1+\varepsilon}$.  For a character trivial on $A$, that
mean is exactly $1$.  The factors $\xi(c_0)$ are $1$ for
retained characters.  Among these, only $1$ is trivial on $A$:
the only other possibility would be the excluded $\chi_q$.
The retained mean therefore gives exactly (CMP.1).  There are
only $4o$ characters to estimate.  Moving $a_t/b$ to $x$ costs
$O(\log x)$ per character by (CMP.12), with the constants in
(CMP.4).  The Fourier bias costs
$O(x^\varepsilon(\log x+1))$, since
$|\mathcal G_\xi(x)|\le x(\log x+C_0+K_0)$ for a retained
nonprincipal character.  All these constants are fixed as the
host grows.  Their sum is at most $4x^{1/2+\varepsilon}$ beyond
one common threshold.  Omitted characters contribute at worst
$-\sum_{\mathcal B\setminus K}|\mathcal G_\xi(x)|$.
With this sum denoted by $\Sigma(x)$, (CMP.13) becomes
\[
 D\{H_\psi(x)-\Sigma(x)-4x^{1/2+\varepsilon}\}
                       \le\overline\Gamma+2Q/x.
\]
By (CMP.4), $\Sigma(x)\le4o x^{1/2+\varepsilon}$.
Divide by $D=Qd_2/2>0$ to obtain (CMP.11).
\end{proof}
% END CMP2

% BEGIN CMP3
\begin{proposition}[CMP3: the mean gap has square-root scale]
\label{claim:CMP3}
If (CMP.9) holds, then for all sufficiently large $x$
\[
                     H_\psi(x)\le x^{1/2+(C+4)\varepsilon}.
 \tag{CMP.14}
\]
\end{proposition}
\begin{proof}
Put $Z=x^{1/2+13\varepsilon/8}$.  Dividing (CMP.7) by $Q$ gives
\[
 \frac BQ=\frac{a_E}{mb}+\frac{a_Em}{bk}
           +\frac{2mx^\varepsilon}{d_1}
           +\frac{d_2kx^{2\varepsilon}}{d_1n^2}.
\]
The four terms are bounded respectively by $2Z,2Z,4Z,2Z$.
Indeed the first is at most $2x/m$, the second at most
$2xm/k\le2$, and the last two use precisely the upper bounds
for $m,k$ and the lower bound for $d_1$ in (CMP.6).
Hence $B\le10QZ$.
Retain $\Sigma$ in the penultimate inequality of CMP2 and divide
by $Q/2$.  By $d_2\ge x^{-\varepsilon}$ and $C\ge0$ this yields
\[
 H_\psi(x)-\Sigma(x)
          \le26x^{1/2+(C+21/8)\varepsilon}
          \le x^{1/2+(C+3)\varepsilon},
\]
the last step holding once $26\le x^{3\varepsilon/8}$.
Now use $\Sigma(x)\le4o x^{1/2+\varepsilon}$ and
$x^\varepsilon\ge1+4o$.  This gives (CMP.14).
The source values $e_D=2,e_L=0$ give
$C=e_D+e_L+3=5$, and hence the exponent $1/2+9\varepsilon$.
This identifies the selected route's constants; their arithmetic
production is still part of the two inputs.
\end{proof}
% END CMP3

% BEGIN CMP4
\begin{theorem}[CMP4: the one-sided analytic criterion]
\label{claim:CMP4}
Fix $0<\theta<1$.  Suppose that for every $\alpha>0$ and every
$Y_0$ there is $Y\ge\max(3,Y_0)$ such that
\[
       \int_Y^{Y^{3/2}}H_\psi(x)_+x^{-\theta-1}\,dx\le Y^\alpha.
 \tag{CMP.15}
\]
Then every constituent of $\mathcal C_\psi$ has no zero in
$\operatorname{Re}s>\theta$, apart from the principal pole.
In particular it is sufficient to have
$H_\psi(x)\le C_\theta x^\theta$ throughout arbitrarily late
windows, with $C_\theta$ fixed.
\end{theorem}
\begin{proof}
Here is the complete pole argument, consolidating AN2--AN3 and AQ3.
For $\operatorname{Re}s>1$ the Mellin transforms are
\begin{align*}
 M_1(s)&=-\frac{\zeta'/\zeta(s)}{s(s-1)}
             -\frac1{(s-1)^2}+\frac{1+\gamma}{s-1},\\
 M_\chi(s)&=-\frac{L'/L(s,\chi)-sL'/L(1,\chi)}{s(s-1)}
                                                    \quad(\chi\ne1).
\end{align*}
The apparent singularities at $1$ are removable.  For the principal
term use $\zeta'/\zeta(s)=-(s-1)^{-1}+\gamma+O(s-1)$.
Set $a_1=1$ and $a_\chi=(3o+1)^{-1}$ otherwise.  The transform
of the real function $H_\psi$ is
\[
 \Phi(s)=\tfrac12\sum_{\chi\in\mathcal C_\psi}
                    a_\chi\{M_\chi(s)+M_{\bar\chi}(s)\}.
\]
It extends meromorphically to $\operatorname{Re}s>0$.  Every
constituent zero $\rho$ there produces a simple pole with residue
\[
 \frac{b_\rho}{\rho(1-\rho)},\qquad
 b_\rho=\tfrac12\sum_\chi a_\chi
                    \{m_\chi(\rho)+m_{\bar\chi}(\rho)\}>0,
                                                               \tag{CMP.16}
\]
where $m_\chi(\rho)$ is its multiplicity, zero if it is not a
zero of that function.  Thus coincident zeros cannot cancel.  A
real pole in $(0,1)$ has positive residue.  The Euler product and
nonvanishing on the line $1$ put all poles in $0<\operatorname{Re}s<1$.

Put $g(u)=H_\psi(e^u)$.  The fixed finite family gives the crude
bound $|g(u)|\le M e^{2u}$ for $u\ge0$.  At
$\sigma=336$, define
\[
 A_k(z)=\frac{(-1)^k}{k!}\Phi^{(k)}(z)
       =\int_0^\infty g(u)\frac{u^k}{k!}e^{-zu}\,du.
\]
For $U=\log Y\ge1$ and $k=\lceil420U\rceil$, the tails obey
\[
 \int_{(0,\infty)\setminus[U,3U/2]}
          \frac{u^k}{k!}e^{-334u}\,du
 \le\frac{e^{-21U}}{320^{k+1}}+
                         \frac{e^{14U}}{348^{k+1}}
 \ll336^{-k}.                                      \tag{CMP.17}
\]
On the upper tail insert $e^{14(u-3U/2)}\ge1$; on the lower
tail insert $e^{14(U-u)}\ge1$, and extend the two integrals to
the positive half-line.  The last inequality uses
$\log(336/320)\le1/20$, $\log(348/336)\ge1/29$,
$420U\le k\le420U+1$.  The constants do not depend on the window.

Reality of $g$ gives, for every real $t$,
\[
 |A_k(\sigma+it)|+A_k(\sigma)
       \le2\int_0^\infty g(u)_+\frac{u^k}{k!}e^{-\sigma u}\,du.
\]
On the good interval, use the elementary bound
$u^ke^{-(\sigma-\theta)u}/k!\le(\sigma-\theta)^{-k}$.
Condition (CMP.15), after $x=e^u$, and (CMP.17) imply
\[
 |A_k(\sigma+it)|+A_k(\sigma)
       \le2e^{\alpha U}(\sigma-\theta)^{-k}+O(\sigma^{-k})
       \le2R_\alpha^{-k}+O(\sigma^{-k}),
 \quad R_\alpha=(\sigma-\theta)e^{-\alpha/420}.
                                                               \tag{CMP.18}
\]
The extra real derivative is what permits a one-sided bound.

Suppose a pole lies to the right of $\theta$, and put
$D_0=\min\{\sigma-\theta,\inf_\rho|\sigma-\rho|\}$.
If some pole has horizontal distance $<D_0$ from
$\operatorname{Re}s=\sigma$, choose $t$ close to its imaginary
part so that $z=\sigma+it$ has a unique nearest pole $\rho$
at distance $a<D_0$.  Only finitely many poles compete in the
relevant compact disks.  Each pair with different imaginary parts
excludes at most one $t$; pairs with the same imaginary part have
different horizontal distances.  Hence such a choice is possible.
Choose $a<r<D_0$ with no other pole in the closed $r$-disk about
$z$.  The $r$-disk about $\sigma$ is pole-free.  Subtraction of
the principal part and Cauchy's derivative estimate give
\[
 A_k(z)=\frac{c}{(z-\rho)^{k+1}}+O(r^{-k}),\qquad
 A_k(\sigma)=O(r^{-k}),\qquad c\ne0.
\]
Choose $\alpha>0$ so small that $R_\alpha>r$, before choosing
arbitrarily late windows.  Equation (CMP.18) would force
$|c|a^{-k-1}\le O(r^{-k})+2R_\alpha^{-k}$ as $k\to\infty$,
which is impossible since $a<r<R_\alpha$.

In the remaining case every pole has horizontal distance at least
$D_0$.  The hypothesized right-hand pole forces
$D_0<\sigma-\theta$.  A nearest pole to $\sigma$ exists in a
compact disk.  It must be the unique real pole
$\beta=\sigma-D_0$, because total distance $D_0$ cannot equal a
larger horizontal distance.  Its residue $c$ is positive by
(CMP.16).  For $D_0<r<\sigma-\theta$ sufficiently close to
$D_0$, Cauchy's estimate gives
$A_k(\sigma)=cD_0^{-k-1}+O(r^{-k})$.
Choose $\alpha$ with $R_\alpha>r$ and apply (CMP.18) at $t=0$.
Its left side is $2A_k(\sigma)>0$ for late $k$, again a
contradiction.  There can therefore be no right-hand pole.

Finally a pointwise bound in the statement makes the integral in
(CMP.15) at most $(C_\theta/2)\log Y$, which is at most
$Y^\alpha$ for sufficiently large $Y$.  This proves that
special case too.
\end{proof}
% END CMP4

% BEGIN CMP5
\begin{theorem}[CMP5: the conditional main conclusion]
\label{claim:CMP5}
If the two arithmetic inputs above hold for every fixed primitive
target with the stated order of choices, then RH and GRH for all
Dirichlet $L$-functions hold.  For RH alone it suffices to supply
the construction for $\psi=1$.
\end{theorem}
\begin{proof}
Fix a proposed constituent zero with real part $\beta>1/2$.
Choose $0<\delta<\min(1/2,\beta-1/2)$ and then
$0<\varepsilon<\min(\varepsilon_0,\delta/(C+4))$.
All thresholds in CMP2--CMP3 are now fixed.  The arithmetic input
supplies arbitrarily late windows beyond them.  At every real
point of each such window, CMP3 gives
$H_\psi(x)\le x^{1/2+\delta}$.
CMP4 at $\theta=1/2+\delta$ contradicts the proposed zero.
The same applies to the principal constituent, whose pole at $1$
has been explicitly removed from the argument.

The primitive functional equation reflects a zero in the open
critical strip to the left of $1/2$ into a zero to the right of
$1/2$ of the conjugate primitive character.  Its gamma, conductor
and root-number factors are nonzero in that strip.  For a character
induced by $\chi^*$,
\[
 L(s,\chi)=L(s,\chi^*)
       \prod_{\substack{p\mid N\\p\nmid\operatorname{cond}\chi^*}}
                                      (1-\chi^*(p)p^{-s}).
\]
The additional factors are nonzero for $\operatorname{Re}s>0$,
since $|\chi^*(p)p^{-s}|<1$.  Thus the assertion passes to all
Dirichlet characters.  The ordinary functional equation and the
known trivial zeros give the usual RH formulation as well.
The scope here is Dirichlet GRH; no assertion about arbitrary
automorphic $L$-functions is inferred.
\end{proof}
% END CMP5

% BEGIN CMP6
\begin{corollary}[CMP6: an even smaller sufficient averaged input]
\label{claim:CMP6}
Fix $1/2<\theta<1$ and choose $0<\varepsilon<\theta-1/2$.
For the same compatible data, suppose that on the selected windows
the nonnegative function
\[
 \mathcal E(x)=\frac{(\overline\Gamma(x))_+}{D(x)}
                                      +\frac4{d_2(x)x}
 \tag{CMP.19}
\]
is measurable and, for every $\alpha>0$, has arbitrarily late
windows on which
$\int_Y^{Y^{3/2}}\mathcal E(x)x^{-\theta-1}\,dx\le Y^\alpha$.
Then the same constituent zero-free half-plane
$\operatorname{Re}s>\theta$ follows, without (CMP.9).
\end{corollary}
\begin{proof}
CMP2 gives
$H_\psi(x)_+\le\mathcal E(x)+(4+4o)x^{1/2+\varepsilon}$.
The weighted integral of the second term is bounded by
$(4+4o)Y^{1/2+\varepsilon-\theta}/
             (\theta-1/2-\varepsilon)$.
Use the hypothesis with $\alpha/2$ and choose the window late
enough that adding this bounded term still leaves at most
$Y^\alpha$.  CMP4 applies.  Measurability is stated because an
arbitrary pointwise choice of auxiliary data need not provide it.
\end{proof}
% END CMP6

\subsection{Where the remaining arithmetic work enters}

CMP1, the cutoff/Fourier part of CMP2, CMP3, CMP4 and CMP5 are
consolidations of existing complete arguments WB1--WB8, SC1,
AN2--AN3 and AQ1--AQ4.  The explicit averaging-before-estimation
in CMP2 and the resulting input (CMP.19) are additional sufficient
weakening results.  Neither constructs the missing arithmetic data.
The new GQT1--GQT5 packet similarly isolates three sufficient row
counts and two bad-set weight estimates for the retained-defect
input; GQT6--GQT7 show that a uniform polynomial spectral error
suffices for the selected global comparison.  These substitutions
may make reconstruction easier, but their actual suppliers remain
to be written on the same family.

For orientation, the source obtains (CMP.9) by comparing integral
period maps with their archimedean norms.  The denominator and
local-gain estimates force divisibility of the same integral
matrices.  The period/height estimate and spectral estimate control
their norms and the size of their common target.  The retained-row
estimate ensures that discarding exceptional rows leaves the
required positive rank.  The arithmetic slope comparison then
produces the budget (CMP.7).  The geometry and deformation theory
construct those matrices, lattices and rows together, rather than
independently.  This explains why the remaining frontier includes
both construction and quantitative statements.

The following table maps every one of the 33 v0.41 leaf identifiers.
The category letters concern the reconstruction: B denotes a
remaining installation, C a remaining human proof of a producer,
and D a remaining exact match to a published theorem.  They are
not claims that the corresponding mathematics is unproved in the
literature.  ``Construction'' refers to the first arithmetic input;
``gap'' to the second on that construction.

\begin{longtable}{p{0.44\linewidth}p{0.06\linewidth}p{0.41\linewidth}}
\textbf{v0.41 leaf} & \textbf{Kind} & \textbf{Role in the conditional proof}\\\hline
\endhead
\texttt{GLOBAL\_DENOMINATOR} & C & Gap: actual local denominator table and bounds.\\
\texttt{GLOBAL\_RETAINED\_DEFECT} & C & Gap: enough actual rows survive with the original ranks.\\
\texttt{GLOBAL\_PERIOD\_HEIGHT} & C & Gap: norms, eigenvalues, metrics and height on the same period maps.\\
\texttt{GLOBAL\_LOCAL\_GAIN} & C & Gap: local divisibility gain with the required uniform loss.\\
\texttt{GLOBAL\_SPECTRAL} & C & Gap: normalized kernel and target-dimension estimates.\\
\texttt{TASK\_PROP\_GK\_INSTALL} & B & Construction: install one-character control and its congruence integer.\\
\texttt{TASK\_ONE\_CHAR\_PROOF} & C & Construction: the actual one-character argument.\\
\texttt{TASK\_INP\_Z\_ASSEMBLY} & B & Construction: the exact Kato input record.\\
\texttt{RIG\_II} & C & Construction: type-II polarized tangent rigidity.\\
\texttt{AD\_VANISH\_K\_MATCH} & D & Construction: full polarized adjoint vanishing; supplied by ADK1--ADK6 here.\\
\texttt{AUT\_UNRAM\_COEF} & C & Construction: automorphy over the actual unramified coefficient field.\\
\texttt{CNR3} & C & Construction: uniform rank-three/rank-two congruence integer.\\
\texttt{COUNT\_II\_SUPPORT} & C & Construction: two compatible rank-four links and same-ideal support.\\
\texttt{TASK\_CASE3\_AUT\_JOIN} & B & Dirichlet construction: exact Case-III automorphy and image join.\\
\texttt{TASK\_CASE3\_PROMOD\_JOIN} & C & Dirichlet construction: original chain-to-promodularity proof.\\
\texttt{SEED\_IV\_ALT} & C & Dirichlet construction: seed with the prechosen residual datum.\\
\texttt{TENSOR\_FACTOR\_TOWER} & B & Dirichlet construction: actual tensor factors and characters.\\
\texttt{CHAIN\_RANK\_FOUR} & C & Dirichlet construction: propagation along the actual chain.\\
\texttt{CHAIN\_DESC4} & B & Dirichlet construction: residual chain descent.\\
\texttt{RANK4\_CHAIN\_REDUCTION} & C & Dirichlet construction: seeded reduction to that chain.\\
\texttt{TASK\_COUNT\_V\_INSTALL} & B & Construction: original finite CountV assembly.\\
\texttt{THM\_S6\_S} & C & Construction: unitary-three ordinary support and automorphy.\\
\texttt{HOST\_FRAME\_LENGTH} & B & Construction: actual host/frame with normalized lengths.\\
\texttt{HOST\_RIGID\_S} & D & Construction: original host's allowed tangents; retained open after ADK6.\\
\texttt{HOST\_NUMBER\_S} & C & Construction: uniform catalogue for the original rank-three hosts.\\
\texttt{TASK\_HOST\_VAN\_E\_INSTALL} & B & Construction: literal HostVanE from the supplied duality.\\
\texttt{TASK\_HOST\_XPP\_INSTALL} & B & Construction: literal HostXpp from the supplied CorX.\\
\texttt{DESIGN\_CHOICES\_COMMON} & C & Common geometry: design family and one original bundle.\\
\texttt{LOCAL\_NBUILT\_COMMON} & C & Common geometry: local specialization on that bundle.\\
\texttt{SECOND\_LAYER\_COMPATIBLE\_SELECTION} & C & Common geometry: compatible second-layer selector.\\
\texttt{BASE\_THETA\_READINGS} & C & Common geometry: actual base readings, maps, metrics and measures.\\
\texttt{BASE\_THETA\_RATIONAL\_TESTS} & C & Common geometry: simultaneous rationality tests for those maps.\\
\texttt{SLR\_BASE\_ROWS\_PRODUCER} & C & Common geometry: actual base row identities.\\
\end{longtable}

This is a map of a selected sufficient route, not a proof that the
route is logically minimal.  In particular the seven Case-III and
chain items are specific to the general Dirichlet extension on this
route; the ordinary target avoids that additional branch.  The
full exact scope of each leaf remains the contract in the v0.41
register, carried into the accompanying JSON without weakening.
ADK6 is the one whole-leaf discharge proposed by this packet.

\subsection{Source appendix for the conditional proof}

All formal-source references are pinned to
\texttt{7cbf7a0c32fca1badf05931879a37c3f8a2b03b4}.
The human proofs reused are
\texttt{window\_bridge/main.tex} (WB1--WB8),
\texttt{analytic\_envelope/analytic\_envelope.tex} (AQ1--AQ5),
and the AN2--AN3 and SC1 sections of
\texttt{RH\_readable\_reconstruction\_v0\_41.tex}.
Their exact hashes, and the source definitions for the original
character, gap and window joins, are recorded in
\texttt{conditional\_main\_v42/contracts.json}.
The correspondence is
\texttt{ChBalance} (CMP1),
\texttt{ChCore}, \texttt{GmAEq}, \texttt{SigALe} (CMP2),
\texttt{CoreGm} (CMP3), and
\texttt{TopWin}, \texttt{landau\_window}, \texttt{lemmaL}
(CMP4--CMP5).  These references are provenance for the
reconstructed implications, not assumptions used to bypass either
unproved arithmetic input.
